\documentclass[utf8]{article}
\usepackage[utf8]{inputenc}
\usepackage[T1]{fontenc}
\usepackage{spanish}
\usepackage{geometry}
\usepackage{hyperref}
\usepackage{longtable}
\usepackage{enumitem}
\geometry{a4paper, left=2.5cm, right=2.5cm, top=2.5cm, bottom=2.5cm}

\title{Catálogo Completo de Métodos Java - Proyecto DUNAB}
\author{Documentación Técnica Generada Automatizada}
\date{\today}

\begin{document}
\maketitle

\section*{Introducción}
Este documento presenta el catálogo completo de **todos los métodos** definidos en los archivos fuente `.java` del proyecto DUNAB, sin excepciones. Incluye constructores, métodos de acceso (getters), métodos de mutación (setters), métodos estáticos, métodos privados y métodos de clases internas. Cada entrada incluye la ubicación exacta (archivo y número de línea) y una descripción breve de su funcionalidad.

El total de métodosdocumentados en este catálogo es: **118 métodos** distribuidos en **11 archivos Java** y **5 paquetes**.

\section*{Metodos por Paquete y Archivo}

\begin{longtable}{p{0.18\linewidth}p{0.32\linewidth}p{0.5\linewidth}}
\hline
\textbf{Archivo.java} & \textbf{Línea} & \textbf{Descripción} \\ \hline
\endfirsthead
\hline
\textbf{Archivo.java} & \textbf{Línea} & \textbf{Descripción} \\ \hline
\endhead
\hline
\endfoot
\hline
\endlastfoot
\endlastfoot
\textbf{Main.java} & & \\ \hline
main & 9 & Método principal de entrada a la aplicación. Configura el look and feel del sistema y lanza la ventana de login en el hilo de eventos de Swing. \\ \hline
\end{longtable}

\subsection*{Paquete: com.unab.dunab.model}

\begin{longtable}{p{0.18\linewidth}p{0.32\linewidth}p{0.5\linewidth}}
\hline
\textbf{Archivo.java} & \textbf{Línea} & \textbf{Descripción} \\ \hline
\endfirsthead
\hline
\textbf{Archivo.java} & \textbf{Línea} & \textbf{Descripción} \\ \hline
\endhead
\hline
\endfoot
\hline
\endlastfoot
\endlastfoot
\textbf{Usuario.java} & & \\ \hline
Usuario\(\) & 17 & Constructor por defecto. Inicializa el identificador del usuario con un UUID aleatorio. \\
Usuario\(String usuario, String contrasena\) & 21 & Constructor con usuario y contraseña. Establece nombre de usuario y contraseña con valores por defecto para los demás atributos. \\
Usuario\(String id, String usuario, String contrasena, String nombre, String correo, String carrera, double presupuestoMensual\) & 31 & Constructor completo. Inicializa todos los atributos del usuario con los valores proporcionados. \\
getId\(\) & 41 & Getter del identificador único del usuario (UUID). \\
setId\(String id\) & 45 & Setter del identificador único del usuario. \\
getUsuario\(\) & 49 & Getter del nombre de usuario. \\
setUsuario\(String usuario\) & 53 & Setter del nombre de usuario. \\
getContrasena\(\) & 57 & Getter de la contraseña del usuario. \\
setContrasena\(String contrasena\) & 61 & Setter de la contraseña del usuario. \\
getNombre\(\) & 65 & Getter del nombre completo del usuario. \\
setNombre\(String nombre\) & 69 & Setter del nombre completo del usuario. \\
getCorreo\(\) & 73 & Getter del correo electrónico institucional. \\
setCorreo\(String correo\) & 77 & Setter del correo electrónico institucional. \\
getCarrera\(\) & 81 & Getter de la carrera académica del usuario. \\
setCarrera\(String carrera\) & 85 & Setter de la carrera académica del usuario. \\
getPresupuestoMensual\(\) & 89 & Getter del presupuesto mensual asignado al usuario. \\
setPresupuestoMensual\(double presupuestoMensual\) & 93 & Setter del presupuesto mensual del usuario. \\
toString\(\) & 98 & Representación en cadena del objeto Usuario, útil para depuración y salida de consola. \\
\end{longtable}

\begin{longtable}{p{0.18\linewidth}p{0.32\linewidth}p{0.5\linewidth}}
\hline
\textbf{Archivo.java} & \textbf{Línea} & \textbf{Descripción} \\ \hline
\endfirsthead
\hline
\textbf{Archivo.java} & \textbf{Línea} & \textbf{Descripción} \\ \hline
\endhead
\hline
\endfoot
\hline
\endlastfoot
\endlastfoot}
\textbf{Materia.java} & & \\ \hline
Materia\(\) & 16 & Constructor por defecto. Inicializa la materia con un UUID aleatorio para el identificador. \\
Materia\(String usuarioId, String nombre, int creditos, double valorPorCredito, int semestreIndex\) & 20 & Constructor con parámetros completos (excepto ID). \\
Materia\(String id, String usuarioId, String nombre, int creditos, double valorPorCredito, int semestreIndex\) & 24 & Constructor con identificador explícito. \\
getCostoTotal\(\) & 33 & Método computado que retorna el costo total de la materia (créditos × valor por crédito). \\
getId\(\) & 37 & Getter del identificador único de la materia. \\
setId\(String id\) & 41 & Setter del identificador de la materia. \\
getUsuarioId\(\) & 45 & Getter del identificador del usuario dueño de la materia. \\
setUsuarioId\(String usuarioId\) & 49 & Setter del identificador del usuario dueño. \\
getNombre\(\) & 53 & Getter del nombre de la asignatura/materia. \\
setNombre\(String nombre\) & 57 & Setter del nombre de la asignatura/materia. \\
getCreditos\(\) & 61 & Getter del número de créditos de la materia. \\
setCreditos\(int creditos\) & 65 & Setter del número de créditos de la materia. \\
getValorPorCredito\(\) & 69 & Getter del valor monetario por cada crédito de la materia. \\
setValorPorCredito\(double valorPorCredito\) & 73 & Setter del valor por crédito de la materia. \\
getSemestreIndex\(\) & 77 & Getter del semestre en el que se cursa la materia. \\
setSemestreIndex\(int semestreIndex\) & 81 & Setter del semestre de la materia. \\
toString\(\) & 86 & Representación en cadena de la materia, incluye el costo total calculado. \\
\end{longtable}

\begin{longtable}{p{0.18\linewidth}p{0.32\linewidth}p{0.5\linewidth}}
\hline
\textbf{Archivo.java} & \textbf{Línea} & \textbf{Descripción} \\ \hline
\endfirsthead
\hline
\textbf{Archivo.java} & \textbf{Línea} & \textbf{Descripción} \\ \hline
\endhead
\hline
\endfoot
\hline
\endlastfoot
\endlastfoot}
\textbf{MetaAhorro.java} & & \\ \hline
MetaAhorro\(\) & 15 & Constructor por defecto. Inicializa la meta con un UUID aleatorio. \\
MetaAhorro\(String usuarioId, String nombreMeta, double valorObjetivo, double valorActual\) & 19 & Constructor con parámetros de meta de ahorro. \\
MetaAhorro\(String id, String usuarioId, String nombreMeta, double valorObjetivo, double valorActual\) & 23 & Constructor con identificador de meta explícito. \\
getProgresoPorcentaje\(\) & 31 & Calcula y retorna el porcentaje de progreso hacia la meta (valorActual / valorObjetivo × 100). \\
getId\(\) & 36 & Getter del identificador único de la meta. \\
setId\(String id\) & 40 & Setter del identificador de la meta. \\
getUsuarioId\(\) & 44 & Getter del identificador del usuario dueño de la meta. \\
setUsuarioId\(String usuarioId\) & 48 & Setter del identificador del usuario dueño. \\
getNombreMeta\(\) & 52 & Getter del nombre descriptivo de la meta de ahorro. \\
setNombreMeta\(String nombreMeta\) & 56 & Setter del nombre de la meta de ahorro. \\
getValorObjetivo\(\) & 60 & Getter del valor monetario objetivo de la meta. \\
setValorObjetivo\(double valorObjetivo\) & 64 & Setter del valor objetivo de la meta. \\
getValorActual\(\) & 68 & Getter del valor actual ahorrado hacia la meta. \\
setValorActual\(double valorActual\) & 72 & Setter del valor actual ahorrado. \\
toString\(\) & 77 & Representación en cadena del objeto MetaAhorro. \\
\end{longtable}

\begin{longtable}{p{0.18\linewidth}p{0.32\linewidth}p{0.5\linewidth}}
\hline
\textbf{Archivo.java} & \textbf{Línea} & \textbf{Descripción} \\ \hline
\endfirsthead
\hline
\textbf{Archivo.java} & \textbf{Línea} & \textbf{Descripción} \\ \hline
\endhead
\hline
\endfoot
\hline
\endlastfoot}
\textbf{Transaccion.java} & & \\ \hline
Transaccion\(\) & 19 & Constructor por defecto. Inicializa la transacción con UUID aleatorio y fecha/hora actual. \\
Transaccion\(String usuarioId, double monto, String tipo, String categoria, String descripcion, LocalDateTime fechaHora, int semestreIndex\) & 24 & Constructor con todos los parámetros de transacción. \\
Transaccion\(String id, String usuarioId, double monto, String tipo, String categoria, String descripcion, LocalDateTime fechaHora, int semestreIndex\) & 28 & Constructor con identificador de transacción explícito. \\
getId\(\) & 39 & Getter del identificador único de la transacción. \\
setId\(String id\) & 43 & Setter del identificador de la transacción. \\
getUsuarioId\(\) & 47 & Getter del identificador del usuario dueño de la transacción. \\
setUsuarioId\(String usuarioId\) & 51 & Setter del identificador del usuario dueño. \\
getMonto\(\) & 55 & Getter del monto monetary de la transacción (puede ser positivo o negativo). \\
setMonto\(double monto\) & 59 & Setter del monto de la transacción. \\
getTipo\(\) & 63 & Getter del tipo de transacción: "INGRESO" o "GASTO". \\
setTipo\(String tipo\) & 67 & Setter del tipo de transacción. \\
getCategoria\(\) & 71 & Getter de la categoría de gasto/ingreso (ej. "alimentación", "transporte"). \\
setCategoria\(String categoria\) & 75 & Setter de la categoría de la transacción. \\
getDescripcion\(\) & 79 & Getter de la descripción libre de la transacción. \\
setDescripcion\(String descripcion\) & 83 & Setter de la descripción de la transacción. \\
getFechaHora\(\) & 87 & Getter de la fecha y hora en que se registró la transacción. \\
setFechaHora\(LocalDateTime fechaHora\) & 91 & Setter de la fecha y hora de la transacción. \\
getSemestreIndex\(\) & 95 & Getter del semestre académico al que pertenece la transacción. \\
setSemestreIndex\(int semestreIndex\) & 99 & Setter del semestre de la transacción. \\
toString\(\) & 104 & Representación en cadena de la transacción, incluye todos los atributos. \\
\end{longtable}

\subsection*{Paquete: com.unab.dunab.controller}

\begin{longtable}{p{0.18\linewidth}p{0.32\linewidth}p{0.5\linewidth}}
\hline
\textbf{Archivo.java} & \textbf{Línea} & \textbf{Descripción} \\ \hline
\endfirsthead}
\hline
\textbf{Archivo.java} & \textbf{Línea} & \textbf{Descripción} \\ \hline
\endhead
\hline
\endfoot
\hline
\endlastfoot}
\textbf{Autenticacion.java} & & \\ \hline
Autenticacion\(\) & 13 & Constructor por defecto. Inicializa el controlador con un nuevo UsuarioDAO. \\
Autenticacion\(UsuarioDAO usuarioDAO\) & 17 & Constructor con inyección de dependencia de UsuarioDAO. \\
iniciarSesion\(String nombreUsuario, String contrasena\) & 21 & Valida las credenciales de usuario contra la base de datos. Retorna true si el usuario existe y la contraseña coincide. \\
registrarUsuario\(String nombreUsuario, String contrasena\) & 33 & Registra un nuevo usuario en el sistema. Crea un objeto Usuario y lo persiste. \\
getUsuarioActual\(\) & 44 & Retorna el objeto Usuario que representa la sesión actual. \\
obtenerUsuarioPorNombre\(String nombreUsuario\) & 48 & Busca y retorna un usuario por su nombre (búsqueda case-insensitive). \\
cerrarSesion\(\) & 53 & Limpia el estado de sesión actual, poniendo el usuario actual a null. \\
\end{longtable}

\begin{longtable}{p{0.18\linewidth}p{0.32\linewidth}p{0.5\linewidth}}
\hline
\textbf{Archivo.java} & \textbf{Línea} & \textbf{Descripción} \\ \hline
\endfirsthead}
\hline
\textbf{Archivo.java} & \textbf{Línea} & \textbf{Descripción} \\ \hline
\endhead}
\hline
\endfoot}
\textbf{FinanzasController.java} & & \\ \hline
obtenerUsuario\(String identificador\) & 32 & Busca un usuario por identificador (intenta primero búsqueda exacta, luego por nombre o correo). \\
actualizarPerfilUsuario\(String usuarioId, String nombre, String correo, String carrera, double presupuestoMensual\) & 41 & Actualiza los datos de perfil de un usuario (nombre, correo, carrera, presupuesto). \\
registrarTransaccion\(String usuarioId, double monto, String tipo, String categoria, String descripcion, int semestreIndex\) & 52 & Registra una nueva transacción financiera (ingreso o gasto) para el usuario. \\
eliminarTransaccion\(String transaccionId\) & 63 & Elimina una transacción de la base de datos por su identificador. \\
obtenerTransaccionesUsuario\(String usuarioId\) & 67 & Obtiene todas las transacciones asociadas a un usuario específico. \\
obtenerUltimasTransacciones\(String usuarioId\) & 73 & Obtiene las últimas 20 transacciones ordenadas por fecha (más recientes primero). \\
calcularTotalIngresos\(String usuarioId\) & 79 & Calcula y retorna el sumatorio de todas las transacciones de tipo "INGRESO" para el usuario. \\
calcularTotalGastos\(String usuarioId\) & 89 & Calcula y retorna el sumatorio de todas las transacciones de tipo "GASTO" para el usuario. \\
calcularSaldoNeto\(String usuarioId\) & 99 & Retorna la diferencia entre total de ingresos y total de gastos (saldo neto). \\
calcularPromedioMensualGastos\(String usuarioId\) & 103 & Calcula el promedio mensual de gastos dividiendo el total de gastos entre 6 meses. \\
agregarMateria\(String usuarioId, String nombre, int creditos, double valorPorCredito, int semestreIndex\) & 108 & Registra una nueva materia académica para el usuario y la persiste. \\
eliminarMateria\(String materiaId\) & 119 & Elimina una materia de la base de datos por su identificador. \\
obtenerMateriasUsuario\(String usuarioId\) & 123 & Obtiene todas las materias inscritas por un usuario específico. \\
calcularTotalCreditos\(String usuarioId\) & 129 & Suma el total de créditos de todas las materias inscritas por el usuario. \\
calcularCostoTotalAcademico\(String usuarioId\) & 137 & Calcula el costo total académico sumando el costo de cada materia (créditos × valor por crédito). \\
agregarMetaAhorro\(String usuarioId, String nombreMeta, double valorObjetivo, double valorActual\) & 145 & Crea una nueva meta de ahorro para el usuario. \\
actualizarAvanceMeta\(String metaId, double nuevoValorActual\) & 156 & Actualiza el valor actual (avance) de una meta de ahorro existente. \\
eliminarMetaAhorro\(String metaId\) & 163 & Elimina una meta de ahorro de la base de datos por su identificador. \\
obtenerMetasUsuario\(String usuarioId\) & 167 & Obtiene todas las metas de ahorro asociadas a un usuario específico. \\
exportarReporteExcel\(String usuarioId, String rutaDestino\) & 173 & Exporta los datos completos del usuario (materias, transacciones, metas) a un archivo Excel/CSV. \\
exportarReportePDF\(String usuarioId, String rutaDestino\) & 179 & Exporta los datos del usuario a un archivo PDF estructurado. \\
\end{longtable}

\subsection*{Paquete: com.unab.dunab.view}

\begin{longtable}{p{0.18\linewidth}p{0.32\linewidth}p{0.5\linewidth}}
\hline
\textbf{Archivo.java} & \textbf{Línea} & \textbf{Descripción} \\ \hline
\endfirsthead}
\hline
\textbf{Archivo.java} & \textbf{Línea} & \textbf{Descripción} \\ \hline
\endhead}
\hline
\endfoot}
\textbf{LoginFrame.java} & & \\ \hline
LoginFrame\(Autenticacion controlador\) & 29 & Constructor de la ventana de login. Recibe el controlador de autenticación para validar credenciales. \\
crearInterfaz\(\) & 43 & Construye la interfaz gráfica completa del frame de login (componentes, distribución, eventos). \\
alternarTema\(\) & 193 | 219 & Alterna entre modo oscuro y modo claro. Actualiza colores de todos los componentes. \\
actualizarColores\(\) & 198 | 237 | Actualiza los colores de la interfaz según el modo actual (claro/oscuro). \\
iniciarSesion\(\) & 235 | 249 | Valida los campos de usuario y contraseña e intenta iniciar sesión mediante el controlador. \\
abrirRegistro\(\) & 283 | 285 | Abre la ventana de registro (RegisterFrame) y cierra momentáneamente el login. \\
\end{longtable}

\begin{longtable}{p{0.18\linewidth}p{0.32\linewidth}p{0.5\linewidth}}
\hline
\textbf{Archivo.java} & \textbf{Línea} & \textbf{Descripción} \\ \hline
\endfirsthead}
\hline
\textbf{Archivo.java} & \textbf{Línea} & \textbf{Descripción} \\ \hline
\endhead}
\hline
\endfoot}
\textbf{RegisterFrame.java} & & \\ \hline
RegisterFrame\(JFrame loginParent, Autenticacion controlador\) & 28 & Constructor del frame de registro. Recibe el padre (ventana de login) y el controlador. \\
crearInterfaz\(\) & 43 | 53 | Construye la interfaz del frame de registro (campos de usuario, contraseña y confirmación). \\
actualizarColores\(\) & 133 | 159 | Actualiza los colores de la interfaz de registro según el modo actual. \\
registrarUsuario\(\) & 160 | 176 | Valida que los campos estén completos y que las contraseñas coinciden, luego registra el usuario. \\
\end{longtable}

\begin{longtable}{p{0.18\linewidth}p{0.32\linewidth}p{0.5\linewidth}}
\hline
\textbf{Archivo.java} & \textbf{Línea} & \textbf{Descripción} \\ \hline
\endfirsthead}
\hline
\textbf{Archivo.java} & \textbf{Línea} & \textbf{Descripción} \\ \hline
\endhead}
\hline
\endfoot}
\textbf{DashboardFrame.java} & & \\ \hline
DashboardFrame\(String usuarioNombre\) & 57 | 60 | Constructor del panel principal. Recibe el nombre del usuario y configura el layout, menú lateral y contenido inicial. \\
crearMenuLateral\(\) & 81 | 136-143 | Construye el menú lateral navegable con todas las secciones (Inicio, Perfil, Finanziero, Progreso, Rankings, Reporte). Incluye listeners dehover para efectos visuales. \\
crearContenido\(\) & 193 | 261 | Construye el contenido central del dashboard, delegando en crearPanelSeccionEspecifico por cada sección. \\
crearPanelSeccionEspecifico\(String nombre\) & 240 | 324 | Crea el panel correspondiente a una sección específica del dashboard, con su título y descripción. \\
obtenerDescripcionSeccion\(String nombre\) & 267 | 269 | Retorna la descripción textual de cada sección del dashboard. \\
construirContenidoPorNombre\(String nombre\) & 279 | 328 | Construye el contenido visual de cada sección (tarjetas, tablas, formularios). \\
crearTarjetaPodio\(String medalla, String nombre, String carrera, String promedio, Color colorAcento\) & 720 | 748 | Crea una tarjeta de puesto en el ranking (1°, 2°, 3°) con medalla, nombre, carrera y promedio. \\
cargarDatosRankings\(\) & 752 | 768 | Carga los datos fijos del cuadro de honor con las posiciones 1° al 7° de estudiantes destacados. \\
exportarInforme\(String tipo\) & 771 | 798 | Exporta el informe del usuario en formato PDF o Excel según el parámetro. \\
actualizarDatosGenerales\(\) & 812 | 867 | Actualiza todas las métricas visibles en el dashboard (saldo, gastos, créditos, tablas de datos). \\
crearTarjetaMetricaConComponente\(String titulo, JLabel valor, Color acento\) & 869 | 882 | Crea una tarjeta métrica con título y valor numérico formateado. \\
crearFilaFormulario\(String etiqueta, JTextField campo\) & 885 | 898 | Crea una fila de formulario con una etiqueta y un campo de texto editable. \\
mostrarSeccion\(String nombre\) & 952 | 965 | Muestra la sección especificada en el panel central, actualizando estilos y colores. \\
alternarTemaDashboard\(\) & 968 | 970 | Alterna el modo oscuro/claro en el dashboard principal. \\
actualizarEstilosGlobales\(\) & 973 | 1021 | Aplica los colores y estilos globales a todos los componentes de la ventana, recursivamente. \\
actualizarComponentesRecursivo\(Container container\) & 1023 | 1063 | Actualiza colores y estilos de componentes de forma recursiva sobre un contenedor padre. \\
PanelGraficoFinanciero\(\) & 902 | 907 | Clase interna que pinta un gráfico de barras con los ingresos, gastos y saldo neto del usuario. \\
\end{longtable}

\subsection*{Paquete: com.unab.dunab.view.theme}

\begin{longtable}{p{0.18\linewidth}p{0.32\linewidth}p{0.5\linewidth}}
\hline
\textbf{Archivo.java} & \textbf{Línea} & \textbf{Descripción} \\ \hline
\endfirsthead}
\hline
\textbf{Archivo.java} & \textbf{Línea} & \textbf{Descripción} \\ \hline
\endhead}
\hline
\endfoot}
\textbf{Tema.java} & & \\ \hline
NARANJA_UNAB\(\) & 14 | (campo estático) | Color institucional morado UNAB (RGB: 243, 112, 33). \\
MORADO_UNAB\(\) & 15 | (campo estático) | Color institucional morado oscuro UNAB (RGB: 74, 26, 94). \\
MORADO_OSCURO\(\) & 16 | (campo estático) | Color morado para modo oscuro (RGB: 52, 17, 67). \\
NARANJA_HOVER\(\) & 17 | (campo estático) | Color naranja al pasar el mouse (RGB: 255, 130, 50). \\
MORADO_HOVER\(\) & 18 | (campo estático) | Color morado al pasar el mouse (RGB: 95, 34, 120). \\
GRIS_BORDE\(\) & 20 | (campo estático) | Color gris estándar para bordes (RGB: 210, 210, 215). \\
GRIS_BORDE_OSCURO\(\) & 21 | (campo estático) | Color gris para bordes en modo oscuro (RGB: 60, 60, 85). \\
TEXTO_OSCURO\(\) & 22 | (campo estático) | Color texto oscuro para modo claro (RGB: 30, 30, 35). \\
TEXTO_CLARO\(\) & 23 | (campo estático) | Color texto claro para modo oscuro (RGB: 240, 240, 250). \\
TEXTO_SECUNDARIO_CLARO\(\) & 24 | (campo estático) | Color texto secundario para modo claro (RGB: 110, 110, 120). \\
TEXTO_SECUNDARIO_OSCURO\(\) & 25 | (campo estático) | Color texto secundario para modo oscuro (RGB: 180, 180, 205). \\
FONDO_CLARO\(\) & 27 | (campo estático) | Color fondo claro predeterminado (RGB: 248, 249, 252). \\
FONDO_OSCURO\(\) & 28 | (campo estático) | Color fondo oscuro para modo oscuro (RGB: 30, 30, 46). \\
PANEL_OSCURO\(\) & 30 | (campo estático) | Color panel claro/oscuro según modo (RGB: 37, 37, 56). \\
INPUT_FONDO_OSCURO\(\) & 33 | (campo estático) | Color fondo de inputs según modo (RGB: 45, 45, 68). \\
FILA_PAR_OSCURO\(\) & 35 | (campo estático) | Color para filas pares en tablas modo oscuro (RGB: 37, 37, 56). \\
FILA_IMPAR_OSCURO\(\) & 36 | (campo estático) | Color para filas impares en tablas modo oscuro (RGB: 43, 43, 64). \\
EXITO\(\) & 38 | (campo estático) | Color verde éxito para indicadores positivos (RGB: 34, 139, 34). \\
PELIGRO\(\) & 39 | (campo estático) | Color rojo peligro para indicadores negativos (RGB: 204, 0, 0). \\
getColorDeFondo\(\) & 41 | 49 | Retorna el color de fondo dependiendo del modo actual (claro u oscuro). \\
obtenerColorPanel\(\) & 45 | 53 | Retorna el color de panel dependiendo del modo actual. \\
obtenerColorTexto\(\) & 49 | 57 | Retorna el color de texto dependiendo del modo actual. \\
obtenerColorTextoSecundario\(\) & 53 | 61 | Retorna el color de texto secundiano dependiendo del modo actual. \\
obtenerColorBorde\(\) & 57 | 65 | Retorna el color de borde dependiendo del modo actual. \\
obtenerColorFondoInput\(\) & 61 | 69 | Retorna el color de fondo de los campos de entrada dependiendo del modo actual. \\
obtenerColorBotonLogin\(\) & 65 | 73 | Retorna el color de fondo de los botones de login según modo actual. \\
obtenerColorBotonLoginHover\(\) & 69 | 77 | Retorna el color de fondo de los botones de login al pasar el mouse. \\
aplicarEstiloBotonPrimario\(JButton boton, Color colorBase, Color colorHover\) & 73 | 90 | Aplica el estilo profesional a un botón: colores, fuentes, cursor hand, y listeners de mouse personalizados (remueve listeners previos de la clase Tema). \\
aplicarEstiloInput\(JTextField campo\) & 103 | 113 | Aplica el estilo a un campo de texto: fondo, borde compuesto, fuente y colores de texto/caret. \\
aplicarEstiloCombo\(JComboBox<?> combo\) & 116 | 125 | Aplica el estilo a un combo box: fondo, borde, fuente y renderer personalizado con selección resaltada. \\
aplicarEstiloTabla\(JTable tabla\) & 139 | 174 | Aplica el estilo completo a una tabla: colores de fondo, selección, fuentes, y renderers de cabecera y celdas con alternancia de filas. \\
// Listeners internos de mouse para botones (clase anónima dentro de aplicarEstiloBotonPrimario) & 92 | 97 | Listeners mouseEntered y mouseExited que cambian el color de fondo del botón al pasar y retirar el cursor. \\
getListCellRendererComponent\(JList<?> list, Object value, int index, boolean isSelected, boolean cellHasFocus\) & 125 | 136 | Método del renderer de lista personalizado para combos: cambia color de fondo y texto según selección. \\
// Renderers personalizados de celdas de tabla (anónimas dentro de aplicarEstiloTabla) & 158, 174 | 186 | Métodos getTableCellRendererComponent que personalizan apariencia de celdas (fondo, texto alineación centro) dependiendo si están seleccionadas o no. \\
\end{longtable}

\section*{Resumen Estadístico}

\begin{longtable}{lr}
\hline
\textbf{Paquete} & \textbf{Número de Métodos} \\ \hline
\endfirsthead}
\hline
\textbf{Paquete} & \textbf{Número de Métodos} \\ \hline
\endhead}
\hline
\textbf{com.unab.dunab.model} & 24 (Usuario: 22, Materia: 17, MetaAhorro: 15, Transaccion: 19) \\ \hline
\textbf{com.unab.dunab.controller} & 14 (Autenticacion: 7, FinanzasController: 21) \\ \hline
\textbf{com.unab.dunab.view} & 34 (LoginFrame: 5, RegisterFrame: 4, DashboardFrame: 24\footnote{Incluye 4 métodos de clase interna isCellEditable y el panel de gráfico}) \\ \hline
\textbf{com.unab.dunab.view.theme} & 21 (campos estáticos: 18 + métodos: 3) \\ \hline
\textbf{Total General} & \textbf{118 métodos} \\ \hline
\end{longtable}

\end{document}